\chapter{清单与模板}\label{s:checklist}

\cite{Gawa2007}普及了「用清单能救命」这个想法，
近期研究总体上支持了清单的有效性~\cite{Avel2013,Urba2014,Rams2019}。
我们觉得它们很有用，
尤其是在把新老师带进团队的时候；
下面给出的这些，可以作为制定你自己清单的起点。

\seclbl{教学评价}{s:checklists-teacheval}

这份评分标准用来评估人们用幻灯片、现场编程或两者混合进行的 5—10 分钟教学。
把每一项评为「是」「勉强」「否」或「不适用」。

\noindent
\begin{longtable}{p{.25\textwidth}p{.75\textwidth}}

  开场
  & 存在（若不存在，其余各项用「不适用」） \\
  & 长度合适（10—30 秒） \\
  & 介绍自己 \\
  & 介绍将要覆盖的主题 \\
  & 说明先修要求 \\
  \\ [-1.5ex] \hline \\ [-1.5ex]

  内容
  & 目标/叙事弧清晰 \\
  & 语言包容 \\
  & 真实任务/例子 \\
  & 教最佳实践/使用地道的代码 \\
  & 在行话的锡拉礁与过度简化的卡律布狄斯之间走出一条航路 \\
  \\ [-1.5ex] \hline \\ [-1.5ex]

  讲授
  & 声音清晰、易懂（口音重则记「勉强」或「否」） \\
  & 节奏：不太快也不太慢，没有长停顿或自我打断，不是明显照稿念 \\
  & 自信：不滑进「不确定」的黏沼、也不掉进「居高临下」的粪堆 \\
  \\ [-1.5ex] \hline \\ [-1.5ex]

  幻灯片
  & 存在（若不存在，其余各项用「不适用」） \\
  & 幻灯片与讲解互补（双重编码） \\
  & 字体和颜色易读/没有铺满整屏的大段文字 \\
  & 屏幕频繁变化（每 30 秒有变化） \\
  & 图形用得好 \\
  \\ [-1.5ex] \hline \\ [-1.5ex]

  现场编程
  & 使用了（若没使用，其余各项用「不适用」） \\
  & 代码与讲解互补 \\
  & 字体和颜色易读/屏幕上代码量合适 \\
  & 工具使用熟练 \\
  & 突出代码的关键特征 \\
  & 剖析错误 \\
  \\ [-1.5ex] \hline \\ [-1.5ex]

  收尾
  & 存在（若不存在，其余各项用「不适用」） \\
  & 长度合适（10—30 秒） \\
  & 总结关键点 \\
  & 勾画下一步 \\
  \\ [-1.5ex] \hline \\ [-1.5ex]

  总体
  & 各点连接清楚/逻辑流畅 \\
  & 让主题有意思（也就是不无聊） \\
  & 有学识 \\

\end{longtable}

\seclbl{团队协作评价}{s:checklists-teameval}

这份评分标准用来评估团队中个人的表现。你可以把它当作创建你自己评分标准的起点。把每一项评为「是」「勉强」「否」或「不适用」。

\noindent
\begin{longtable}{p{.25\textwidth}p{.75\textwidth}}

  沟通
  & 专注倾听他人，不打断 \\
  & 复述他人所说以确保理解 \\
  & 清楚、简洁地表达想法 \\
  & 为想法给出好理由 \\
  & 赢得他人支持 \\
  \\ [-1.5ex] \hline \\ [-1.5ex]

  决策
  & 从不同角度分析问题 \\
  & 在解决问题时运用逻辑 \\
  & 基于事实而非「直觉」或感觉提出方案 \\
  & 向他人征求新想法 \\
  & 产生新想法 \\
  & 接受变化 \\
  \\ [-1.5ex] \hline \\ [-1.5ex]

  协作
  & 承认团队需要面对和解决的问题 \\
  & 朝着所有相关方都能接受的方案努力 \\
  & 与他人分享成功的功劳 \\
  & 鼓励所有参与者参与 \\
  & 公开、不设防地接受批评 \\
  & 与他人合作 \\
  \\ [-1.5ex] \hline \\ [-1.5ex]

  自我管理
  & 监控进展以确保目标达成 \\
  & 把取得结果放在最优先 \\
  & 为工作会议确定任务优先级 \\
  & 鼓励他人表达哪怕是相反的观点 \\
  & 会议期间专注于任务 \\
  & 高效利用会议时间 \\
  & 在工作会议期间提出推进方式 \\

\end{longtable}

\seclbl{活动筹备}{s:checklists-events}

下面的清单用于活动之前、之中和之后。

\subsection*{安排活动}

\begin{itemize}

\item
  决定是线下、
  单站点在线，
  还是多站点在线。

\item
  和主办方把预期谈清楚，
  确保大家都同意由谁承担差旅费。

\item
  确定谁可以参加：
  活动是对所有人开放、
  仅限某个组织的成员，
  还是介于两者之间？

\item
  安排教师。

\item
  安排场地，需要的话包括分组讨论室。

\item
  选定日期。
  如果是线下，就订好行程。

\item
  向主办方获取参加者的姓名和邮箱地址。

\item
  确保每个人都完成注册。

\end{itemize}

\subsection*{搭建}

\begin{itemize}

\item
  建一个网页，写明工作坊详情，
  包括日期、
  地点、
  以及参与者需要带什么。

\item
  检查是否有参加者有特殊需求。

\item
  如果是线上工作坊，测试视频会议——测两次。

\item
  确保参加者能上网。

\item
  建一个用于共享笔记和练习解答的页面（例如一个 Google Doc）。

\item
  给参加者发一封欢迎邮件，里面带上
  工作坊页面的链接、
  背景阅读材料、
  对他们需要做的配置的说明、
  需要携带物品的清单、
  以及当天联系主办方或教师的方式。

\end{itemize}

\subsection*{活动开始时}

\begin{itemize}

\item
  提醒大家行为准则。

\item
  点名，
  并整理一份姓名清单，粘贴到共享页面里。

\item
  分发便利贴。

\item
  确保每个人都能上网。
  
\item
  确保每个人都能访问共享页面。

\item
  收集任何相关的在线账号 ID。

\end{itemize}

\subsection*{活动结束时}

\begin{itemize}

\item
  更新出席名单。

\item
  向参与者收集反馈。

\item
  把共享页面复制一份。

\end{itemize}

\subsection*{差旅包}

下面是一些老师会带去工作坊的东西：

\begin{longtable}{p{0.45\textwidth}p{0.45\textwidth}}

便利贴 & 润喉糖 \\
舒服的鞋 & 一个小记事本 \\
备用电源适配器 & 备用衬衫 \\
除臭剂 & 视频转接头 \\
笔记本电脑贴纸 & 他们的笔记（打印件或平板上的） \\
一根燕麦棒或其他零食 & 抗酸药（因为路上吃的） \\
名片 & 备用眼镜/隐形眼镜 \\
一个笔记本和笔 & 一支激光笔 \\
用于茶/咖啡的保温杯 & 额外的白板笔 \\
牙刷或漱口水 & 湿巾（因为总会洒出来） \\

\end{longtable}

出行时，
许多老师还会带上跑鞋、泳衣、瑜伽垫，
或者随便什么他们运动时穿的、用的东西。
有些人还会带一个便携式 WiFi 热点，以防房间的网络不工作，
以及几个装着学习者所需软件安装程序的 U 盘。

\seclbl{课程设计}{s:checklists-design}

本节总结由~\cite{Wigg2005,Bigg2011,Fink2013}各自独立提出的逆向设计方法。
它给出一个分步推进的流程，
帮你按正确的顺序思考正确的事情，
并提供有间隔的交付物，
好让你能重新界定范围或调整方向，而不会遇到太多令人不快的意外。

从第 2 步往后的所有东西都会进入你的最终课程，
所以不会有白费的功夫；
正如\chapref{s:process}里描述的，
早点写示例练习有助于确保
你要求学习者做的每件事都对课程目标有贡献，
而且他们需要知道的东西都被覆盖到了。

这些步骤是按细节递增的顺序描述的，
但过程本身始终是迭代的。
当你从后面问题的答案里学到东西、
或意识到最初的计划不会像你当初想的那样展开时，
你经常会回过头去修改前面的工作。

\subsection*{这门课是给谁的？}

创建一些学习者画像（\secref{s:process-personas}），
或者（最好是）挑一些你和同事为通用目的画好的。
每个人物画像应该有：

\begin{enumerate}

\item
  此人的大致背景；

\item
  他已经会什么；

\item
  他认为自己想做什么；以及

\item
  他有什么特殊需求。

\end{enumerate}

~\\
\noindent
\textbf{交付物：}你试图帮助的是谁的简要描述。

\subsection*{大想法是什么？}

用要点形式回答下面三四个问题，
帮你弄清这门课的范围。
你不必回答所有这些问题，
如果觉得有用，也可以自己提出并回答别的，
但第一个问题你总该给出几条答案。
在这个阶段你还可以画一张概念图（\secref{s:memory-concept-maps}）。

\begin{itemize}

\item
  人们将学会解决什么问题？

\item
  人们将学到哪些概念和技巧？

\item
  人们将使用哪些技术、包或函数？

\item
  你将定义哪些术语或行话？

\item
  你将用哪些类比来解释概念？

\item
  你预期会有哪些错误或误解？

\item
  你将使用哪些数据集？

\end{itemize}

~\\
\noindent
\textbf{交付物：}
课程的大致范围。
把它分享给一位同事——这时候的一点点反馈，
能省下日后好几小时的无效劳动。

\subsection*{学习者在过程中会做什么？}

通过写出学习者在这门课接近尾声时能够做的几道练习的完整描述，
把第 2 步的目标定得更实。
这么做类似\hreffoot{https://en.wikipedia.org/wiki/Test-driven\_development}{测试驱动开发}：
不是从一组（可能含糊的）学习目标向前推，
而是从你想让学习者最终达到的具体例子往回推。
这还有助于发现技术上的一些要求，
否则这些要求可能要晚到令人难受的时候才会暴露出来。

为了配合这些完整练习描述，
每一讲课时写一两道练习的要点式描述，
以显示你预期学习者的进展有多快。
再说一遍，
这些是对「你假设了多少东西」的良好现实检验，
也有助于发现技术要求。
创建这些「额外」练习的一种办法，
是把解决主要练习所需的技能列成要点清单，
再针对每一项各出一道练习。

~\\
\noindent
\textbf{交付物：}1—2 道完整讲解的练习，
用上人们要学的技能，
外加六道要点式练习提纲。
附上完整的解答，
好让你确认为学习者准备的软件真的能跑。

\subsection*{概念之间如何连接？}

把你创建的练习排成一个合理的顺序，
再由此推导出一份要点式课程提纲。
提纲应当每小时有三四个要点，
每个都配上某种形成性评价。
在这个阶段改动评价很常见，
好让它们能相互支撑。

~\\
\noindent
\textbf{交付物：}一份课程提纲。
在这个阶段，你很可能会发现以前忘了列出的东西，
所以如果你得反复回头几次，别惊讶。

\subsection*{课程概览}

现在你可以写一份课程概览，包括：

\begin{itemize}

\item
  一段描述（也就是给学习者的推销词）；

\item
  六七个学习目标；以及

\item
  先修要求的总结。

\end{itemize}

提前做这件事往往白费功夫，
因为材料在更早的步骤里通常会被增删、挪动。

~\\
\noindent
\textbf{交付物：}
课程描述、
学习目标、
和先修要求。

\seclbl{课前评估问卷}{s:checklists-preassess}

这份问卷帮教师估量入门 JavaScript 工作坊参与者的既有编程知识。
问题和答案都很具体，
而且整份很短，好让作答者不会被吓到。

\begin{enumerate}

\item
  下面哪一项最能描述
  你总体上编程的经验？

  \begin{itemize}
    
  \item
    我没有。
    
  \item
    我偶尔写过几行。
    
  \item
    我为自用写过几页长的程序。
    
  \item
    我编写并维护过更大的软件。\\
    
  \end{itemize}

\item
  下面哪一项最能描述
  你用 JavaScript 编程的经验？

  \begin{itemize}
    
  \item
    我没有。
    
  \item
    我偶尔写过几行。
    
  \item
    我为自用写过几页长的程序。
    
  \item
    我编写并维护过更大的软件。\\
    
  \end{itemize}

\item
  下面哪一项最能描述
  你有多容易用任何语言写一个程序，
  找出一个列表里最大的数？

  \begin{itemize}
    
  \item
    我不知道从哪儿下手。
    
  \item
    我靠大量网上搜索、反复试错，能磕磕绊绊做出来。
    
  \item
    我几乎不用外部帮助就能很快做出来。\\
    
  \end{itemize}

\item
  下面哪一项最能描述
  你有多容易写一个 JavaScript 程序，
  找出一个网页里所有的标题并把它们首字母大写？

  \begin{itemize}
    
  \item
    我不知道从哪儿下手。
    
  \item
    我靠大量网上搜索、反复试错，能磕磕绊绊做出来。
    
  \item
    我几乎不用外部帮助就能很快做出来。\\
    
  \end{itemize}

\item
  上完这门课之后，你想知道什么、或者能做到什么，
  是你现在不知道、或做不到的？

\end{enumerate}
